% Generated draft: requires mathematical and source review.
\documentclass[11pt]{article}
\usepackage[margin=25mm]{geometry}
\usepackage{fontspec}
\setmainfont{Latin Modern Roman}
\usepackage{amsmath,amssymb,mathtools,unicode-math}
\setmathfont{Latin Modern Math}
\usepackage{enumitem,needspace,xurl,hyperref}
\setlength{\parindent}{0pt}
\setlength{\parskip}{5pt}
\setlength{\emergencystretch}{3em}
\newcommand{\N}{\mathbb{N}}
\newcommand{\Z}{\mathbb{Z}}
\newcommand{\Q}{\mathbb{Q}}
\newcommand{\R}{\mathbb{R}}
\newcommand{\C}{\mathbb{C}}
\newcommand{\F}{\mathbb{F}}
\newcommand{\A}{\mathbb{A}}
\newcommand{\PP}{\mathbb{P}}
\newcommand{\E}{\mathbb{E}}
\newcommand{\eps}{\varepsilon}
\newcommand{\dd}{\,\mathrm{d}}
\begin{document}
\section*{UnsolvedMath ID: 20002084}
\textbf{Source problem code:} catalogue entry
\subsection{Definitions and mathematical statement}
Question 1.26 records Eliashberg's proposal to use symplectic field theory (SFT) or ideas inspired by it for the Audin conjecture and weak boundary rigidity of the Clifford torus in \(S^3\). It asks how much needs full SFT foundations, proposes replacing boundary conditions on a Lagrangian by punctured-curve asymptotics in \(T^*L\), and points to a desired study of \(T^*T^n\). The CR-structure suggestion is not formulated precisely. The references to unfinished foundations and to the problem being open are historical 2003 workshop wording, not current claims.

The aid reports the Audin torus statement as solved by Cieliebak--Mohnke: an embedded Lagrangian torus in \(\C^n\) has minimal Maslov number two. It attributes this to targeted neck stretching and punctured-curve transversality, rather than construction of the entire SFT algebra. Exact hypotheses and the primary citation are not supplied; no broader claim about every closed Lagrangian topology is inferred from the torus formulation.

The aid's period reduction is explicit. On \(\C^n\) use
\[
 \omega_0=\sum_i\dd x_i\wedge\dd y_i,\qquad
 \lambda_0=\tfrac12\sum_i(x_i\dd y_i-y_i\dd x_i),
 \quad\dd\lambda_0=\omega_0.
\]
For a closed Lagrangian embedding \(\iota_t:L\to\C^n\), call the isotopy exact when
\(\beta_t=\iota_t^*(\iota_{\partial_t\iota_t}\omega_0)\) is exact for every \(t\). Then Cartan's formula implies
\[
 [\iota_t^*\lambda_0]=[\iota_0^*\lambda_0]\in H^1(L;\R).
\]
Consequently the entire relative-disc area homomorphism, transported using the domain identification, is preserved:
\[
 \pi_2(\C^n,\iota_tL)\longrightarrow\R,
 \qquad[u]\longmapsto\int_u\omega_0
       =\int_{\partial u}\lambda_0.
\]
An exact isotopy does not mean that the individual compact Lagrangians have exact Liouville forms. Hamiltonian isotopies are exact in this deformation sense.

Let
\[
 \Lambda_n=\{|z_1|=\cdots=|z_n|=1/\sqrt n\}
 \subset S^{2n-1}=\partial\overline B^{2n}(1).
\]
Its period group is \((\pi/n)\Z\). Thus a torus exactly isotopic to \(\Lambda_n\) has the same least positive disc area \(\pi/n\). Equivalently, using the infimum of positive areas when needed, its value is still \(\pi/n\); the discrete period group here avoids any ambiguity of a minimum. A torus in the geometric closed unit ball with this value is called extremal in the aid's normalization.

In dimension four, the aid reports Dimitroglou Rizell's published extremal-torus theorem: such a torus in \(\overline B^4(1)\) lies in \(S^3\). Therefore a torus Hamiltonian isotopic in \(\C^2\) to \(\Lambda_2\) and contained in that closed ball must lie in its boundary. Its area group is \((\pi/2)\Z\). This proves the reported weak boundary-rigidity conclusion; it does not assert that the final torus equals the standard torus as a subset, nor require the whole ambient isotopy to preserve the ball. It also precludes such a final torus strictly inside the open unit ball.

The same formal reduction would prove the analogous statement in all dimensions if the corresponding extremal-torus boundary theorem holds. The aid ties that further step specifically to Faisal's revised 2026 preprint, conditionally, not as an independently verified published theorem.

Finally, the broader proposed full-SFT computation for \(T^*T^n\) is not verified by the aid. Rigorous wrapped-Floer/loop-space results supply more limited cotangent structures and do not automatically compute the full SFT algebra, all punctured-curve operations or their foundations. No definite theorem about CR structures is supplied. The terminal 8 and ``Fourth day'' are layout transitions, not additional mathematical conditions.
\subsection{Short English statement}
The aid separates targeted punctured-curve results from a full SFT construction. Exact isotopy preserves the Liouville period group, making a Clifford-isotopic torus in the unit ball extremal. The reported four-dimensional extremal theorem gives weak boundary rigidity; higher-dimensional extension remains tied to the stated conditional preprint.
\subsection{Sources}
\begin{itemize}
\item Permanent UnsolvedMath ID \(20002084\), exact public code catalogue entry; local registry \texttt{aim-geometry-notes.json}, \texttt{notes[421]}, canonical tag question.
\item AIM workshop \emph{Holomorphic curves in contact geometry}, version dated 15 October 2003, Question 1.26, recording Eliashberg's discussion: \url{https://aimath.org/WWN/contactgeom2/contactgeom2.pdf}. No PDF was accessed. The supplied verification note announces recoverable OCR/layout errors but does not list them in the field.
\item The aid attributes the Audin result to Cieliebak--Mohnke, the published four-dimensional extremal-torus result to Dimitroglou Rizell, and a conditional higher-dimensional theorem to Faisal's revised 2026 preprint. Full titles, theorem statements and publication/preprint identifiers are absent. Wrapped-Floer/loop-space work is mentioned without specific citations.
\item Research attempt 1, ``Targeted punctured curves versus full SFT, with a Liouville-period bridge to Clifford boundary rigidity,'' labelled partial\_result, partially\_solved, proved and low novelty confidence. Original literature and proposer/year metadata are absent beyond the supplied source history.
\end{itemize}
\end{document}
